\chapter{Abstract}
Kubernetes does not impose a specific network plugin, forcing teams to select a \emph{Container Network Interface} (CNI) without objective comparative metrics. This situation leads to infrastructure oversizing (CPU/RAM) and incomplete security postures. This thesis builds reproducible empirical evidence —infrastructure as code, automated benchmarks, and a multi-criteria decision analysis (MCDA) model— to guide this selection. Flannel, Calico, Cilium, and Antrea are evaluated on a high-availability K3s cluster on DigitalOcean, measuring TCP throughput, establishment latency, retransmissions, and CPU/RAM overhead of the CNI agent, as well as the impact of Zero Trust network policies. The results feed a parameterizable MCDA model by use profile (Fintech/security, Streaming/performance, and IoT/limited resources) and a web prototype that displays the recommendations.
The main findings show that no plugin proved superior across all criteria: Calico achieved the highest TCP throughput (1804.40 Mbps) and the lowest connection-establishment latency, while Flannel recorded the lowest retransmission rate and resource footprint. Flannel —the default CNI in K3s— was also found to accept NetworkPolicies through the Kubernetes API without enforcing them in its data plane, constituting a security false positive, and Cilium was the only plugin capable of enforcing layer-7 (HTTP) policies. Consistent with this absence of dominance, the model yields a different recommendation for each profile: Cilium for Fintech, Calico for Streaming, and Flannel for IoT, the latter by a narrow margin.

\vspace{1cm}
\textbf{Keywords:} Kubernetes, CNI, NetworkPolicy, Zero Trust, MCDA, QoS, eBPF.
\newpage